% Statement and cited-input packet. Target proof is separately embargoed.
% Frozen II source SHA256: 74497bfe327fae25c9cdcb17334b3bf42a9793a9df7c0a4404d97568f2f760f0
% Target owner: cdp:st:E16; tag: 01GN

\begin{proposition}[Rank divisors from the root Kodaira--Spencer morphism]\hypertarget{tag-01GN}{}\label{cdp:st:E16}


We use the logarithmic Kodaira--Spencer morphism at the specified degenerate fibres and the ordinary morphism on the finite smooth torus covers. For the original affine-twist family, including every $|p|>1$, its only rank defect is rank one at the order-three point. The logarithmic cusp has rank two. For the $(3k,4l)$ logarithmic transforms, on their full root covers the ranks at the two multiple values are respectively
\[
 \begin{cases}1&k=1,\\0&k>1,\end{cases}
 \qquad
 \begin{cases}2&l=1,\\0&l>1.\end{cases}                 \tag{E.16.1}\hypertarget{tag-01A2}{}\label{cdp:eq:E.16.1}
\]
Every other smooth point and the logarithmic cusp have rank two.

For the two-parameter family the logarithmic rank is one at each of its two rank-one nodal values, zero at the order-four root point, and two everywhere else, including the hexagonal fibre. Its determinant has simple coarse zeros at the two nodal values and order four in the root coordinate at the multiple value. For the Liu--Rollenske central family the logarithmic rank is two at both the cuspidal and hexagonal fibres and throughout the remaining smooth locus; at the order-four root point it is zero, and its determinant has order four there.

For a specified positive comb section $\sigma$ of degree parameter $m$, let $x,y$ be root coordinates of weights $3,4$, with $y=0$ at the order-three point. There is a nonzero homogeneous polynomial $P_\sigma(X,Y)$ of degree $4m$, uniquely determined up to an overall nonzero scalar by the marked period section, such that
\[
 \det\kappa_\sigma=y^2P_\sigma(x^4,y^3).             \tag{E.16.2}\hypertarget{tag-01A3}{}\label{cdp:eq:E.16.2}
\]
The rank-defect set is exactly its zero divisor. At any such point the rank is zero exactly when all four entries of the actual Kodaira--Spencer matrix vanish; otherwise it is one. Formula \eqref{cdp:eq:E.16.3} below computes those entries from $\sigma$. No unspecified choice of a comb section is assigned a fixed root set. The polynomial is nonzero at the hexagonal cusp, so that cusp always has rank two.
\end{proposition}

% Cited input cdp:st:E13; statement or labelled formula only.
\begin{proposition}[The Kuranishi germ of the two-parameter family]\hypertarget{tag-01GK}{}\label{cdp:st:E13}


For every member of the family of
\cite[Proposition labelled \texttt{family}]{Viaclovsky2026Two}, on the parameter
bidiscs with $\mu$ sufficiently close to zero and nonzero sufficiently
small linearization scale $\lambda$, one has
\[
 (h^0,h^1,h^2,h^3)(X_{\mu,\lambda},T_X)=(1,2,2,0).
\]
Its Kuranishi germ is a smooth reduced surface, and the displayed
$(\mu,\lambda)$-family is locally miniversal. In particular its
sufficiently small deformations all have algebraic dimension one.
\end{proposition}

% Cited input cdp:st:E14; statement or labelled formula only.
\begin{proposition}[The Liu--Rollenske family and its transverse deformation]\hypertarget{tag-01GL}{}\label{cdp:st:E14}


Every member $Y(u)$ of the Liu--Rollenske family, with its scale
sufficiently small and nonzero as in its construction, satisfies
\[
 h^\bullet(Y(u),T_{Y(u)})=(1,2,2,0).
\]
Its Kuranishi germ is a smooth reduced surface. One parameter is
the given scale; a second parameter replaces its cuspidal fibre by
two rank-one nodal fibres. All sufficiently small deformations
retain a torus fibration and have algebraic dimension one.
\end{proposition}

% Cited input cdp:st:E15; statement or labelled formula only.
\begin{proposition}[Tangent cohomology and Kuranishi germs of positive combs]\hypertarget{tag-01GM}{}\label{cdp:st:E15}


For each positive-comb manifold $X_m$ constructed in
\cite{Viaclovsky2026Combs}, in particular for every $m\geq12$ in
its stated construction, one has
\[
 h^\bullet(X_m,T_{X_m})=(2m+1,\,8m+1,\,6m+1,\,0).
\]
Its Kuranishi germ is smooth and reduced of dimension $8m+1$.
A locally miniversal family is supplied by all sufficiently nearby
admissible period sections on the root curve, with the affine
attachment classes fixed. Its members have algebraic dimension one
and torus fibrations with the same three critical values,
multiplicities $3,4$, and hexagonal fibre type.
\end{proposition}

% Cited input cdp:st:E7; statement or labelled formula only.
\begin{proposition}[Application to the named six-sphere families]\hypertarget{tag-01GW}{}\label{cdp:st:E7}


Every member of the following torus-fibred families,
including the period family $X_c$, its members with $|p|>1$,
Viaclovsky's $X_{3k,4l}$, his two-parameter family and his combs,
and the Liu--Rollenske family, satisfies the rank hypothesis of \ref{cdp:st:E6}.
Consequently every sufficiently small deformation retains a
holomorphic torus fibration over $\mathbb P^1$. Every such member
with $b_2=0$, and every sufficiently small deformation of it, has
algebraic dimension one. In particular the fibration-retention
question in Section~10(1) of the supplied comb paper has an
affirmative answer.
\end{proposition}

% Cited input cdp:st:E8; statement or labelled formula only.
\begin{proposition}[The theta determinant and its exact rank locus for $X_c$]\hypertarget{tag-01GX}{}\label{cdp:st:E8}


Use the period convention
$\Omega=\left(\begin{smallmatrix}\tau&\mu\\\mu&\beta/6\end{smallmatrix}\right)$,
Jacobi's odd theta function with periods $1,\tau$, and
$\beta/6=(i/\pi)\log(\vartheta_1(\mu,\tau)/\eta(\tau))+\text{constant}$.
Writing $h=\partial_\mu\log\vartheta_1$, one has the identity
\begin{equation}\hypertarget{tag-01A1}{}\label{cdp:eq:E-theta-general}
 \det\frac{d\Omega}{d\tau}
 =-\left(\mu'-\frac{i}{2\pi}h\right)^2
       -\frac{\wp(\mu,\tau)}{4\pi^2}.
\end{equation}
For the actual AC period branch $\wp(\mu,\tau)=0$, this becomes
\[
 \det\frac{d\Omega}{d\tau}
       =-\frac{g_2(\tau)^2}{36\pi^2g_3(\tau)}.
\]
On the smooth root-family base, the Kodaira--Spencer tensor has rank
one at the order-three point and rank two at every other finite
point, including the order-four point. Its determinant has order
two at the order-three point. The logarithmic tensor at the cusp
has rank two.
\end{proposition}

% Cited input cdp:st:E6; statement or labelled formula only.
\begin{proposition}[Stability of a rank-two torus fibration]\hypertarget{tag-01GV}{}\label{cdp:st:E6}


Let $g:Y\to\mathbb P^1$ be as in \ref{cdp:st:E5} and suppose its
Kodaira--Spencer tensor has rank two at one regular value. For every
proper deformation $\mathcal Y\to(S,0)$ of $Y$ over a complex
analytic germ, including a nonreduced germ, there is, after shrinking
$S$, a holomorphic map
\[
 G:\mathcal Y\longrightarrow\mathbb P^1\times S
\]
over $S$ whose central restriction is $g$. Every $G_s$ is surjective
with connected fibres and smooth general fibre a complex two-torus.
With $b_2(Y)=0$, every sufficiently small deformation has algebraic
dimension one.
\end{proposition}

% Cited input cdp:st:E5; statement or labelled formula only.
\begin{proposition}[The infinitesimal map and the rank-two boundary]\hypertarget{tag-01GU}{}\label{cdp:st:E5}


Let $g:Y\to C$ be a proper surjective holomorphic map with connected fibres
from a smooth connected compact complex threefold to a nonsingular compact
connected complex curve. Its smooth general fibre is a two-dimensional complex
torus. There is a canonical exact sequence of vector spaces
\begin{equation}\hypertarget{tag-019Y}{}\label{cdp:eq:E-Leray-map}
0\to H^1(C,T_C)\xrightarrow{g^*}H^1(Y,g^*T_C)
 \to H^0(C,T_C\otimes R^1g_*\mathcal O_Y)\to0.
\end{equation}
On the smooth locus $U\subset C$, the relative tangent sequence
induces a sheaf morphism
\[
 \delta:T_C\otimes R^1g_*\mathcal O_Y|_U
                  \longrightarrow R^2g_*T_{Y/C}|_U.
\]
At a smooth torus fibre $F$, write its Kodaira--Spencer tensor as
$\kappa=u_1\otimes\eta_1+u_2\otimes\eta_2$, where
$\eta_1,\eta_2$ is a basis of $H^1(F,\mathcal O_F)$ and
$u_i\in\operatorname{Lie}(F)$. Then
\[
 \delta(a\eta_1+b\eta_2)
   =(bu_1-au_2)\otimes(\eta_1\wedge\eta_2).
\]
Its rank is the rank of $\kappa$. When this rank is two at one
regular value, the image of $H^1(Y,T_Y)\to H^1(Y,g^*T_C)$ lies
in $g^*H^1(C,T_C)$. For $C=\mathbb P^1$ this image is zero.
\end{proposition}

